\section{Conclusion}
\label{sec:conclusion}

We introduce Corpus-Level Inconsistency Detection (\task), addressing the challenge of identifying contradictory information within large knowledge repositories. To tackle this problem, we present \system, an agent-based system that combines retrieval with LLM reasoning to detect and contextualize potential contradictions for human review.

We also release \dataset, a benchmark capturing real inconsistencies that synthetic datasets often miss. Our experiments show that, while retrieval and verification are challenging, \system enables humans to uncover substantially more inconsistencies. Applied to Wikipedia, this framework reveals that approximately 3.3\% of facts conflict with other information in the corpus, amounting to millions of contradictory statements across the encyclopedia.

These results demonstrate that corpus-level inconsistencies are a measurable phenomenon in large-scale knowledge corpora. Although automated systems still exhibit systematic errors, they can aid in maintaining knowledge consistency at scale. More broadly, this work suggests a virtuous cycle: LLMs help curate cleaner, more reliable corpora, which in turn improve both human knowledge access and the AI systems built on top of them.
